\documentclass[UTF8,12pt]{ctexart}
\usepackage[a4paper,margin=24mm]{geometry}
\usepackage{amsmath,amssymb,bm,graphicx,adjustbox}
\newcommand{\fitmath}[1]{\begin{adjustbox}{max width=\linewidth}$\displaystyle #1$\end{adjustbox}}
\setlength{\parindent}{0pt}
\setlength{\parskip}{.7em}
\sloppy
\allowdisplaybreaks
\begin{document}
\section*{1991 年考研数学三 · 参考答案}
\subsection*{1991年真题参考答案}

\subsection*{（试卷IV）}

\subsection*{一、填空题}

（1）\(\displaystyle e^{\sin(xy)}\cos(xy)(y\,dx+x\,dy)\)。 （2）\(\displaystyle -1;-1;1\)。 （3）\(\displaystyle -(n+1);-e^{-(n+1)}\)。


（4）


\[\fitmath{\displaystyle \begin{pmatrix}O&B^{-1}\\A^{-1}&O\end{pmatrix}.}\]


（5）


\[\fitmath{\displaystyle \begin{array}{c|ccc}X&-1&1&3\\\hline P&0.4&0.4&0.2\end{array}}\]


\subsection*{二、选择题}

（1）A。 （2）D。 （3）B。 （4）D。 （5）B。


三、\(\displaystyle e^{(n+1)/2}\)。


四、\(\displaystyle \dfrac{\displaystyle ab^2}{\displaystyle 30}\)。


五、\(\displaystyle y=x\sqrt{2(\ln x+1)}\;(x\ge e^{-1})\)。可求得 \(\displaystyle y^2=2x^2(\ln|x|+1)\)。但是，满足初始条件 \(\displaystyle y|_{x=e}=2e\) 的特解为 \(\displaystyle y=x\sqrt{2(\ln x+1)}\;(x\ge e^{-1})\)。其余三个连续函数 \(\displaystyle y=-x\sqrt{2(\ln x+1)}\;(x\ge e^{-1})\)、\(\displaystyle y=-x\sqrt{2[\ln(-x)+1]}\;(x\le-e^{-1})\)、\(\displaystyle y=x\sqrt{2[\ln(-x)+1]}\;(x\le-e^{-1})\) 都不满足 \(\displaystyle y|_{x=e}=2e\)。


六、\(\displaystyle a=3\)。


七、当 \(\displaystyle p_1=80,\allowbreak p_2=120\) 时，总利润最大，最大利润为605。


八、证明略（\(\displaystyle f(x)=e^{x\ln(1+1/x)}\)，计算 \(\displaystyle f\prime(x)\)，并证明 \(\displaystyle f\prime(x)>0\)）。


九、（1）当 \(\displaystyle \lambda\ne0\) 且 \(\displaystyle \lambda\ne-3\) 时，\(\displaystyle \boldsymbol\beta\) 可由 \(\displaystyle \boldsymbol\alpha_1,\allowbreak \boldsymbol\alpha_2,\allowbreak \boldsymbol\alpha_3\) 唯一地线性表示。（2）当 \(\displaystyle \lambda=0\) 时，表达式不唯一。（3）当 \(\displaystyle \lambda=-3\) 时，不能线性表示。


十、\(\displaystyle -2<\lambda<1\)。


十一、证明略（记 \(\displaystyle A=(\boldsymbol\alpha_1,\allowbreak \ldots,\allowbreak \boldsymbol\alpha_n)\)，线性无关的充分必要条件是 \(\displaystyle |A|\ne0\)，利用 \(\displaystyle |A^{\mathrm T}A|=|A|^2=D\)）。


十二、\(\displaystyle P\{X=0\}=\dfrac{\displaystyle 1}{\displaystyle 2},\allowbreak \ P\{X=1\}=\dfrac{\displaystyle 1}{\displaystyle 4},\allowbreak \ P\{X=2\}=\dfrac{\displaystyle 1}{\displaystyle 8},\allowbreak \ P\{X=3\}=\dfrac{\displaystyle 1}{\displaystyle 8}\)。


十三、（1）\(\displaystyle \rho=0\)。（2）\(\displaystyle X\) 和 \(\displaystyle Y\) 不独立。


十四、\(\displaystyle \hat\lambda=\dfrac{\displaystyle n}{\displaystyle \sum_{i=1}^nX_i^a}\)。


\subsection*{（试卷V）}

\subsection*{一、填空题}

（1）【同试卷IV 第一、（1）题】 （2）【同试卷IV 第一、（2）题】 （3）【同试卷IV 第一、（3）题】


（4）\(\displaystyle a^n+(-1)^{n+1}b^n\)。 （5）\(\displaystyle 0.6\)。


\subsection*{（试卷V）参考答案续}

\subsection*{二、选择题}

（1）【同试卷IV 第二、（1）题】 （2）D。 （3）C。 （4）D。 （5）【同试卷IV 第二、（4）题】


三、\(\displaystyle 1\)。 四、\(\displaystyle \dfrac{\displaystyle 22}{\displaystyle 3}\)。


五、\(\displaystyle x\arctan x-\dfrac{\displaystyle 1}{\displaystyle 2}\ln(1+x^2)-\dfrac{\displaystyle 1}{\displaystyle 2}(\arctan x)^2+C\)，其中 \(\displaystyle C\) 为任意常数。


六、证明略（可对 \(\displaystyle xy=xf(z)+yg(z)\) 两边同时分别关于 \(\displaystyle x,\allowbreak y\) 求偏导数）。


七、【同试卷IV 第六题】 八、【同试卷IV 第七题】


九、证明略（可考虑函数 \(\displaystyle f(x)=\ln(1+\dfrac{\displaystyle 1}{\displaystyle x})-\dfrac{\displaystyle 1}{\displaystyle 1+x}\)，计算 \(\displaystyle f\prime(x)\) 并利用 \(\displaystyle f(x)\) 的单调性）。


十、（1）证明略（注意到 \(\displaystyle AB-A-B+E=(A-E)(B-E)=E\)）。


（2）


\[\fitmath{\displaystyle A=\begin{pmatrix}1&\dfrac{\displaystyle 1}{\displaystyle 2}&0\\-\dfrac{\displaystyle 1}{\displaystyle 3}&1&0\\0&0&2\end{pmatrix}.}\]


十一、【同试卷IV 第九题】


十二、\(\displaystyle k=1\) 或 \(\displaystyle k=-2\)。


十三、（1）【同试卷IV 第十二题】 （2）\(\displaystyle \dfrac{\displaystyle 67}{\displaystyle 96}\)。


十四、（1）\(\displaystyle \alpha=0.064\)。 （2）\(\displaystyle \beta\approx0.009\)。

\end{document}
